\section{总结与展望}

\subsection{完成工作}
本设计按计划交付了一个可演示的云边端协同 AI 手表原型，功能清单如下：
\begin{itemize}
  \item F407 端：LVGL 六行主页与各子页、16\,kHz/16\,bit 录音与外部 SRAM 缓冲、
        MPU6050 计步（阈值可现场改）与跌倒粗判、DS18B20 表温、RTC 与闹钟、
        秒表和倒计时、计算器、三款小游戏、EEPROM 持久化，以及 PE5 点阵滚动；
  \item C3 边缘：Wi-Fi 管理与在线上报、帧协议对端实现、PCM 流式封装 WAV 上传
        火山方舟语音模型、Open-Meteo 天气摘要、文本回传；
  \item 系统级：自定义可靠帧协议（CRC16 + 停等重传 + 链路质量统计）、
        断网降级、版本哈希自检、主机侧仿真测试集。
\end{itemize}


\subsection{收获与体会}
\begin{enumerate}[label=（\arabic*）]
\item 云边端协同的本质是\textbf{把不可靠的东西集中到一层}。网络依赖全部收敛在 C3 及以上，
F407 的软件因此保持了嵌入式应有的确定性。联调中语音链路出问题时，计步与时钟仍然稳定。
\item 自检机制的价值高于事后调试。PCM 写后读校验、CRC/ORE/DROP 计数、
SHA-256 版本横幅，各自在第一时间拦住了一类真实故障。
\item 复用冲突必须前置核对。摄像头、TF 卡、RS232、以太网等与本设计引脚冲突的资源，
全部在设计期用一张冲突表消掉，运行期没有再出现复用引发的故障。
\end{enumerate}

\subsection{展望}
\begin{enumerate}[label=（\arabic*）]
\item \textbf{音频压缩}。引入 IMA-ADPCM，把数据量压到约四分之一后再上云，
缩短 10\,s 录音的串口传输时间。
\item \textbf{端侧关键词唤醒}。在 F407 或 C3 上做离线关键词识别，
减少每次交互都要联网的等待，与云端大模型互补。
\item \textbf{跌倒上报}。启用已定义的 0x30 告警帧，把告警推送到云端或手机。
\item \textbf{产品化}。做 PCB 集成、可佩戴外形和低功耗，并考虑心率、血氧等传感。
\end{enumerate}
